\section{Shortest Routes I}
\label{sec:cses-shortest-routes-i}

\problemstatement{Given a directed graph of $n$ cities and $m$ roads with
non-negative lengths, output the shortest distance from city $1$ to every
city.}

\begin{patternbox}
Single source, non-negative weights: \textbf{Dijkstra's algorithm}. Always
finalise the unvisited node with the smallest tentative distance (found via
a min-heap), then relax its outgoing edges.
\end{patternbox}

\subsection*{Greedy finalisation with a heap}

Keep $\textit{dist}[1]=0$ and all others $\infty$. Repeatedly pop the
smallest $(\textit{d},u)$ from a priority queue; if it is stale (a better
distance is already recorded) skip it, else relax each edge $u\to v$ of
weight $w$: if $\textit{dist}[u]+w<\textit{dist}[v]$, update and push $v$.
Non-negativity guarantees a popped node's distance is final.

\begin{ideabox}[Non-Negativity Makes the Nearest Node Final]
When the closest unfinished node is popped, no later path can beat its
distance, because any such path would pass through a node already at least
as far and then add a non-negative edge. This is why Dijkstra is correct
only without negative edges.
\end{ideabox}

\subsection*{Algorithm}

\begin{enumerate}
  \item $\textit{dist}[1]=0$; push $(0,1)$.
  \item Pop smallest; skip if stale; else relax all out-edges, pushing
        improved neighbours.
  \item Output $\textit{dist}[1\ldots n]$.
\end{enumerate}

\subsection*{Dry run}

\dryrun{Edges $1\!\to\!2(5)$, $1\!\to\!3(2)$, $3\!\to\!2(1)$.}

\begin{itemize}
  \item Pop $(0,1)$: relax $2\to5$, $3\to2$. Pop $(2,3)$: relax $2$ via
        $3$: $2+1=3<5$.
  \item Distances: $\textit{dist}=[0,3,2]$.
\end{itemize}

\subsection*{C++ implementation}

\begin{lstlisting}
#include <bits/stdc++.h>
using namespace std;

int main() {
    int n, m;
    cin >> n >> m;
    vector<vector<pair<int,long long>>> adj(n + 1);   // {to, weight}
    for (int i = 0; i < m; i++) {
        int a, b; long long w; cin >> a >> b >> w;
        adj[a].push_back({b, w});
    }

    vector<long long> dist(n + 1, LLONG_MAX);
    dist[1] = 0;
    priority_queue<pair<long long,int>, vector<pair<long long,int>>,
                   greater<>> pq;
    pq.push({0, 1});
    while (!pq.empty()) {
        auto [d, u] = pq.top(); pq.pop();
        if (d > dist[u]) continue;                     // stale
        for (auto& [v, w] : adj[u])
            if (dist[u] + w < dist[v]) {
                dist[v] = dist[u] + w;
                pq.push({dist[v], v});
            }
    }
    for (int i = 1; i <= n; i++) cout << dist[i] << " ";
    cout << "\n";
    return 0;
}
\end{lstlisting}

\begin{complexitybox}
\textbf{Time Complexity.} $O((n+m)\log n)$. \textbf{Space Complexity.}
$O(n+m)$.
\end{complexitybox}

\begin{takeawaybox}
Shortest Routes I is textbook Dijkstra: heap of tentative distances, skip
stale pops, relax outgoing edges. It is the single-source backbone the
Flight Discount and Flight Routes twists build on.
\end{takeawaybox}
